%----------------------------------------------------------------------------------------
% SYMMETRIEËN EN PERMUTATIES. Alle opgaven zijn eigen werk.
%
% Opbouw: Start, leerstof (theorie, werkwijze of uitleg, voorbeeld, korte oefening), en op een
% nieuwe pagina de Opdrachten (door elkaar, stijgende moeilijkheid), Koppel en Test jezelf.
% Labels van dit hoofdstuk beginnen met h3-.
%----------------------------------------------------------------------------------------

\mychapter{3}{Symmetrieën en permutaties}

\lkr{Leerplan B+S'': LPD 43, met de twee eindige voorbeelden die het leerplan noemt: symmetriegroepen van figuren met de samenstelling (de gelijkzijdige driehoek) en permutaties met de samenstelling. Sectie 3 vergelijkt groepen met de orde van hun elementen. Het woord \textit{isomorf} staat alleen in de uitleg: morfismen horen bij het complementair gedeelte. Hou kartonnen driehoeken en vierkanten met genummerde hoekpunten klaar: leerlingen die de figuur echt draaien en omklappen, maken veel minder fouten bij het samenstellen.}

Op de voorkant van dit boek staat een behangpatroon van William Morris. Schuif je het patroon een stuk op, of spiegel je het om een verticale lijn, dan valt het precies op zichzelf. Zulke bewegingen heten symmetrieën. Een sneeuwvlok heeft er twaalf, een vierkant acht, en een kristal in een mineraal heel wat meer.

In dit hoofdstuk zie je dat de symmetrieën van een figuur een groep vormen, met als bewerking: de ene beweging na de andere uitvoeren. Het zijn je eerste groepen die niet commutatief zijn. Daarna bekijk je permutaties, het door elkaar schudden van een rij voorwerpen. Galois ontdekte de groepen precies zo: door de oplossingen van een vergelijking van plaats te wisselen. Op het einde vergelijk je groepen met elkaar, en zie je wanneer twee groepen eigenlijk dezelfde zijn.

\begin{oefeningen}[Start]
\oef{h3-start-1}

\oef{h3-verken}
\end{oefeningen}

\lkr{Laat oefening \ref{oef:h3-verken} maken vóór de uitleg, met een kartonnen driehoek waarvan de hoekpunten op beide kanten genummerd zijn. In vraag c vinden de leerlingen zelf dat de volgorde van twee bewegingen een verschil maakt.}

%========================================================================================
\section{Symmetrieën van een figuur}

\subsection{De symmetriegroep}

Een spiegeling, een draaiing en een verschuiving veranderen de plaats van een figuur, maar niet haar vorm of grootte: alle afstanden blijven gelijk. Brengt zo'n beweging een figuur precies op zichzelf, dan is ze een symmetrie van die figuur.

\begin{theorie}
Een \textbf{\textcolor{accent}{isometrie}} is een transformatie van het vlak die alle afstanden bewaart, zoals een spiegeling, een draaiing of een verschuiving.

\vspace{1mm}
Een \textbf{\textcolor{accent}{symmetrie}} van een figuur $F$ is een isometrie die $F$ op zichzelf afbeeldt.

\vspace{1mm}
De \textbf{\textcolor{accent}{samenstelling}} $\formule{s \circ t}$ (lees: $s$ na $t$) is de transformatie die eerst $t$ en daarna $s$ uitvoert.

\begin{stelling}{Symmetriegroep}
De symmetrieën van een figuur vormen met de samenstelling een groep: de \textbf{\textcolor{accent}{symmetriegroep}} van de figuur.
\end{stelling}

\begin{bewijs}
Beelden $s$ en $t$ de figuur $F$ op zichzelf af, dan ook $s \circ t$, en $s \circ t$ bewaart afstanden: $\circ$ is een bewerking. Samenstellen is associatief, want voor elk punt $P$ is
\bewijsregel{\bigl((s \circ t) \circ u\bigr)(P) = s\bigl(t(u(P))\bigr) = \bigl(s \circ (t \circ u)\bigr)(P) .}
Het neutraal element is de identieke transformatie $e$, die elk punt op zijn plaats laat. De inverse van $s$ is de beweging die $s$ ongedaan maakt: ook die bewaart afstanden en beeldt $F$ op $F$ af.
\end{bewijs}
\end{theorie}

Let op de volgorde: in $s \circ t$ voer je eerst $t$ uit, de rechtse, net zoals bij samengestelde functies. Omdat de bewerking een samenstelling is, schrijven we de inverse van $s$ als $s^{-1}$, en $s^2$ voor $s \circ s$.

\subsection{De gelijkzijdige driehoek}

In de Verken vond je zes manieren om een kartonnen driehoek terug in zijn omtrek te leggen. Je kunt hem op zijn plaats laten ($e$), draaien over $120^\circ$ of $240^\circ$ tegen de wijzers van de klok in ($r$ en $r^2$), of omklappen om een van de drie symmetrieassen ($s_1$, $s_2$ en $s_3$). De assen liggen vast in het vlak: $s_1$ is de spiegeling om de as door de bovenste hoek, $s_2$ om de as door de hoek linksonder, en $s_3$ om de as door de hoek rechtsonder. Ze bewegen niet mee met de driehoek.

\figuren{\begin{tikzpicture}[scale=.75,font=\small]
  \driehoek[0,0]{1}{2}{3}\node at (0,-1.35) {$e$};
  \driehoek[3.2,0]{3}{1}{2}\node at (3.2,-1.35) {$r$};\draaipijl[3.2,0]{.42}{}
  \driehoek[6.4,0]{2}{3}{1}\node at (6.4,-1.35) {$r^2$};
  \driehoek[9.6,0]{1}{3}{2}\node at (9.6,-1.35) {$s_1$};\draw[grijs,dashed] (9.6,1.15) -- (9.6,-.6);
  \driehoek[12.8,0]{3}{2}{1}\node at (12.8,-1.35) {$s_2$};\draw[grijs,dashed] ($(12.8,0)+(210:1.15)$) -- ($(12.8,0)+(30:.6)$);
  \driehoek[16,0]{2}{1}{3}\node at (16,-1.35) {$s_3$};\draw[grijs,dashed] ($(16,0)+(330:1.15)$) -- ($(16,0)+(150:.6)$);
\end{tikzpicture}}

Een symmetrie ligt vast als je weet waar de hoekpunten terechtkomen. Daarom nummer je de hoekpunten: de figuur toont de driehoek na elk van de zes symmetrieën, vertrekkend van de stand $e$.

\voorbeeld{}{Bereken $s_1 \circ r$ en $r \circ s_1$.\label{vb:h3-samenstellen}}{%
\uitleg{We vertrekken telkens van de beginstand, voeren eerst de rechtse symmetrie uit en dan de linkse. Het resultaat vergelijken we met de zes standen hierboven.}
\figuren{\begin{tikzpicture}[scale=.7,font=\small]
  \driehoek[0,0]{1}{2}{3}
  \draw[-{Stealth[length=2mm]},thick] (1.15,.1) -- node[above] {$r$} (2.15,.1);
  \driehoek[3.3,0]{3}{1}{2}
  \draw[-{Stealth[length=2mm]},thick] (4.45,.1) -- node[above] {$s_1$} (5.45,.1);
  \driehoek[6.6,0]{3}{2}{1}
  \node[anchor=west] at (7.8,.1) {$= s_2$};
  \begin{scope}[xshift=10.6cm]
  \driehoek[0,0]{1}{2}{3}
  \draw[-{Stealth[length=2mm]},thick] (1.15,.1) -- node[above] {$s_1$} (2.15,.1);
  \driehoek[3.3,0]{1}{3}{2}
  \draw[-{Stealth[length=2mm]},thick] (4.45,.1) -- node[above] {$r$} (5.45,.1);
  \driehoek[6.6,0]{2}{1}{3}
  \node[anchor=west] at (7.8,.1) {$= s_3$};
  \end{scope}
\end{tikzpicture}}
\begin{stappenbreed}
s_1 \circ r = s_2 & Eerst $r$: hoekpunt 3 komt boven. Dan $s_1$: de hoeken onderaan wisselen. Hoekpunt 2 blijft linksonder staan, zoals bij $s_2$. \\[1ex]
r \circ s_1 = s_3 & Eerst $s_1$: 2 en 3 wisselen. Dan $r$: alles schuift een plaats op. Hoekpunt 3 staat rechtsonder, zoals bij $s_3$.
\end{stappenbreed}
De uitkomsten verschillen: de symmetriegroep van de driehoek is niet commutatief.}

Zo kun je alle 36 samenstellingen uitrekenen. Dat geeft de tabel van de symmetriegroep van de gelijkzijdige driehoek. In de rij van $a$ en de kolom van $b$ staat $a \circ b$: eerst $b$, dan $a$.

\begin{center}
\begin{bewtabel}[7.5mm]{6}{\circ}{e & r & r^2 & s_1 & s_2 & s_3}
e & e & r & r^2 & s_1 & s_2 & s_3 \\ \hline
r & r & r^2 & e & s_3 & s_1 & s_2 \\ \hline
r^2 & r^2 & e & r & s_2 & s_3 & s_1 \\ \hline
s_1 & s_1 & s_2 & s_3 & e & r & r^2 \\ \hline
s_2 & s_2 & s_3 & s_1 & r^2 & e & r \\ \hline
s_3 & s_3 & s_1 & s_2 & r & r^2 & e \\ \hline
\end{bewtabel}
\end{center}

Lees er een paar dingen uit af. Elke spiegeling is haar eigen inverse, en $r^{-1} = r^2$: wie $120^\circ$ draait en dan $240^\circ$, is rond. Twee draaiingen na elkaar geven een draaiing, twee spiegelingen na elkaar ook. De tabel is niet symmetrisch om de diagonaal: met zes elementen is dit de kleinste groep die niet commutatief is.

\begin{werkwijze}
Teken de figuur met genummerde hoekpunten in de beginstand. Voer eerst de rechtse symmetrie uit en daarna de linkse, en zoek de eindstand in de lijst van symmetrieën. Spiegelassen en draaicentrum liggen vast in het vlak: ze draaien niet mee met de figuur.
\end{werkwijze}

\begin{oefeningen}
\oef{h3-driehoek}
\end{oefeningen}

\subsection{Andere figuren}

Elke figuur heeft haar eigen symmetriegroep. Hoe minder regelmatig de figuur, hoe kleiner de groep.

\voorbeeld{}{Bepaal de symmetriegroep van een rechthoek die geen vierkant is, en stel de tabel op.\label{vb:h3-rechthoek}}{%
\uitleg{We zoeken eerst alle symmetrieën: draaiingen en spiegelingen die de rechthoek op zichzelf leggen. Daarna stellen we samen met de genummerde hoekpunten.}
\figuren{\begin{tikzpicture}[scale=.62,font=\small]
  \rechthoek[0,0]{1}{2}{3}{4}\node at (0,-1.4) {$e$};
  \rechthoek[3.6,0]{3}{4}{1}{2}\node at (3.6,-1.4) {$d$};\fill[grijs] (3.6,0) circle (1.5pt);
  \rechthoek[7.2,0]{4}{3}{2}{1}\node at (7.2,-1.4) {$h$};\draw[grijs,dashed] (5.8,0) -- (8.6,0);
  \rechthoek[10.8,0]{2}{1}{4}{3}\node at (10.8,-1.4) {$v$};\draw[grijs,dashed] (10.8,-.95) -- (10.8,.95);
\end{tikzpicture}}
\naast{\begin{minipage}{\linewidth}
\begin{stappenvrij}{1.8cm}{7.0cm}
e, d & Laten staan, en de halve draai $d$ over $180^\circ$. Een kwartdraai legt de lange zijde verticaal: geen symmetrie. \\[1ex]
h, v & De spiegelingen om de horizontale en de verticale as. De diagonalen zijn geen symmetrieassen. \\[1ex]
h \circ v = d & Eerst $v$, dan $h$: de stand van $d$.
\end{stappenvrij}

\vspace{1mm}
Zo vind je de hele tabel. De groep is commutatief en elk element is zijn eigen inverse.
\end{minipage}}{\begin{bewtabel}{4}{\circ}{e & d & h & v}
e & e & d & h & v \\ \hline
d & d & e & v & h \\ \hline
h & h & v & e & d \\ \hline
v & v & h & d & e \\ \hline
\end{bewtabel}}}

Een regelmatige veelhoek met $n$ zijden heeft $2n$ symmetrieën: $n$ draaiingen om het middelpunt (over een veelvoud van $\frac{360^\circ}{n}$, ook $e$) en $n$ spiegelingen. Een vierkant heeft er dus 8, een regelmatige vijfhoek 10. Er zijn ook figuren met alleen draaiingen, zoals de driearmige figuur hieronder: drie draaiingen, geen spiegeling. Een cirkel heeft oneindig veel symmetrieën, en een fries zoals de rand van de voorkant ook: je kunt hem over één motief verschuiven, of over twee, of terug.

\figuren{\begin{tikzpicture}[font=\small,scale=.92]
  \begin{scope}
  \foreach \h in {45,90,135,180}{\draw[grijs,dashed,thin] (\h:1.15) -- (\h+180:1.15);}
  \draw[accent,thick,fill=accent!8] (.75,.75) -- (-.75,.75) -- (-.75,-.75) -- (.75,-.75) -- cycle;
  \node at (0,-1.55) {8 symmetrieën};
  \end{scope}
  \begin{scope}[xshift=3.6cm]
  \foreach \k in {0,...,4}{\draw[grijs,dashed,thin] (90+72*\k:1.15) -- (270+72*\k:.95);}
  \draw[accent,thick,fill=accent!8] (90:1) -- (162:1) -- (234:1) -- (306:1) -- (18:1) -- cycle;
  \node at (0,-1.55) {10 symmetrieën};
  \end{scope}
  \begin{scope}[xshift=7.2cm]
  \foreach \h in {90,210,330}{\draw[accent,thick,fill=accent!8,rotate=\h,line join=round] (0,0) -- (1.05,0) -- (.7,.5) -- cycle;}
  \node at (0,-1.55) {3 symmetrieën};
  \end{scope}
  \begin{scope}[xshift=10.2cm]
  \foreach \k in {0,...,4}{\draw[accent,thick,fill=accent!8] (\k*.9,-.35) -- (\k*.9+.6,0) -- (\k*.9,.35) -- (\k*.9+.2,0) -- cycle;}
  \node at (-.4,0) {$\cdots$};\node at (4.6,0) {$\cdots$};
  \draw[-{Stealth[length=2mm]},fel!85!black,thick] (0,.7) -- node[above] {$t$} (.9,.7);
  \node at (2.1,-1.55) {oneindig veel};
  \end{scope}
\end{tikzpicture}}

\begin{oefeningen}
\oef{h3-figuren}
\end{oefeningen}

%========================================================================================
\section{Permutaties}

\subsection{Permutaties samenstellen}

Elke symmetrie van de driehoek schudt de drie hoekpunten door elkaar. Laat je de figuur weg en kijk je alleen naar het door elkaar schudden, dan krijg je een permutatie. In de cursus Telproblemen was een permutatie een rangschikking. Hier bekijken we ze als een functie: waar gaat elk element naartoe?

\begin{theorie}
Een \textbf{\textcolor{accent}{permutatie}} van $\{1, 2, \ldots, n\}$ is een functie $\sigma$ van $\{1, 2, \ldots, n\}$ naar zichzelf die elk getal precies één keer als beeld heeft. We noteren ze in twee rijen, met onder elk getal zijn beeld:

\vspace{1mm}
\qquad $\formule{\sigma = \perm{1 & 2 & \cdots & n \\ \sigma(1) & \sigma(2) & \cdots & \sigma(n)}}$

\vspace{1mm}
De verzameling van alle permutaties van $\{1, 2, \ldots, n\}$ noteren we $\formule{S_n}$.

\begin{stelling}{Symmetrische groep}
$(S_n, \circ)$ is een groep met \quad $\formule{|S_n| = n!}$ \quad elementen.
\end{stelling}

\begin{bewijs}
Stel $\sigma$ en $\tau$ samen: verschillende getallen blijven verschillend, dus $\tau \circ \sigma$ is weer een permutatie. Het samenstellen van functies is associatief. Het neutraal element is de permutatie $e$ die elk getal op zijn plaats laat. De inverse $\sigma^{-1}$ stuurt elk beeld $\sigma(k)$ terug naar $k$, en is ook een permutatie.

Voor $\sigma(1)$ zijn er $n$ keuzes, voor $\sigma(2)$ nog $n - 1$, en zo verder tot 1 keuze voor $\sigma(n)$. Volgens de productregel zijn er $n(n - 1) \cdots 1 = n!$ permutaties.
\end{bewijs}
\end{theorie}

\naast{\begin{minipage}{11.2cm}
$S_n$ heet de \textit{symmetrische groep}. Ze groeit heel snel: $S_3$ heeft 6 elementen, $S_5$ al 120, en $S_{10}$ meer dan drie miljoen.

\vspace{1mm}
Een permutatie kun je ook tekenen als pijlen van de bovenste rij naar de onderste: de figuur toont $\sigma = \perm{1 & 2 & 3 & 4 \\ 3 & 1 & 4 & 2}$. Uit elk getal vertrekt één pijl, en bij elk getal komt één pijl aan.
\end{minipage}}{\begin{tikzpicture}[font=\small]
  \pijlen{3,1,4,2}
  \node[anchor=east] at (.3,.5) {$\sigma$};
\end{tikzpicture}}

\voorbeeld{}{Neem $\sigma = \perm{1 & 2 & 3 & 4 \\ 3 & 1 & 4 & 2}$ en $\tau = \perm{1 & 2 & 3 & 4 \\ 2 & 4 & 3 & 1}$. Bereken $\tau \circ \sigma$ en $\sigma \circ \tau$.\label{vb:h3-perm}}{%
\uitleg{Bij $\tau \circ \sigma$ werkt eerst $\sigma$. We volgen elk getal: eerst naar zijn beeld in $\sigma$, dan van daar naar zijn beeld in $\tau$. In de figuur staan de pijlen van $\sigma$ boven die van $\tau$.}
\naast{\begin{minipage}{10.6cm}
\begin{stappenvrij}{4.3cm}{5.8cm}
1 \xrightarrow{\sigma} 3 \xrightarrow{\tau} 3, \quad 2 \xrightarrow{\sigma} 1 \xrightarrow{\tau} 2 & Volg 1 en 2 door beide permutaties. \\[1ex]
3 \xrightarrow{\sigma} 4 \xrightarrow{\tau} 1, \quad 4 \xrightarrow{\sigma} 2 \xrightarrow{\tau} 4 & En 3 en 4. \\[1ex]
\tau \circ \sigma = \perm{1 & 2 & 3 & 4 \\ 3 & 2 & 1 & 4} & De beelden onder elkaar. \\[3ex]
\sigma \circ \tau = \perm{1 & 2 & 3 & 4 \\ 1 & 2 & 4 & 3} & Op dezelfde manier, maar eerst $\tau$: bijvoorbeeld $3 \xrightarrow{\tau} 3 \xrightarrow{\sigma} 4$.
\end{stappenvrij}
\end{minipage}}{\begin{tikzpicture}[font=\small]
  \begin{scope}[yshift=1cm]\pijlen{3,1,4,2}\end{scope}
  \pijlen{2,4,3,1}
  \node[anchor=east] at (.3,1.5) {$\sigma$};
  \node[anchor=east] at (.3,.5) {$\tau$};
\end{tikzpicture}}
Ook $S_4$ is dus niet commutatief.}

\subsection{De inverse permutatie}

De inverse maakt alles ongedaan: ze stuurt elk beeld terug naar waar het vandaan kwam. In de notatie met twee rijen draai je dus de pijlen om.

\voorbeeld{}{Bepaal de inverse van $\sigma = \perm{1 & 2 & 3 & 4 & 5 \\ 4 & 5 & 1 & 3 & 2}$.\label{vb:h3-perminvers}}{%
\uitleg{De inverse stuurt 4 terug naar 1, 5 naar 2, enzovoort. We wisselen de twee rijen, en zetten de kolommen daarna in de gewone volgorde.}
\begin{stappenbreed}
\perm{4 & 5 & 1 & 3 & 2 \\ 1 & 2 & 3 & 4 & 5} & Wissel de bovenste en de onderste rij. \\[3ex]
\sigma^{-1} = \perm{1 & 2 & 3 & 4 & 5 \\ 3 & 5 & 4 & 1 & 2} & Sorteer de kolommen volgens de bovenste rij.
\end{stappenbreed}
Controle: $1 \xrightarrow{\sigma} 4 \xrightarrow{\sigma^{-1}} 1$ en $2 \xrightarrow{\sigma} 5 \xrightarrow{\sigma^{-1}} 2$, en zo verder: $\sigma^{-1} \circ \sigma = e$.}

\begin{werkwijze}
Voor $\tau \circ \sigma$ volg je elk getal $k$: eerst naar $\sigma(k)$ in de rechtse permutatie, dan naar $\tau\bigl(\sigma(k)\bigr)$ in de linkse. Voor $\sigma^{-1}$ wissel je de twee rijen van $\sigma$ en zet je de kolommen in volgorde. Een vergelijking $\sigma \circ x = \tau$ los je op zoals in elke groep: $x = \sigma^{-1} \circ \tau$.
\end{werkwijze}

\nujij{Bereken $\sigma \circ \sigma$ en $\sigma^{-1}$ voor $\sigma = \perm{1 & 2 & 3 \\ 3 & 1 & 2}$.}{$1 \to 3 \to 2$, $2 \to 1 \to 3$, $3 \to 2 \to 1$, dus $\sigma \circ \sigma = \perm{1 & 2 & 3 \\ 2 & 3 & 1}$. Rijen wisselen en sorteren: $\sigma^{-1} = \perm{1 & 2 & 3 \\ 2 & 3 & 1}$. Dat is hetzelfde, want $\sigma^3 = e$.}

\begin{oefeningen}
\oef{h3-perm}
\end{oefeningen}

%========================================================================================
\section{Dezelfde structuur}

\subsection{De orde van een element}

De symmetriegroep van de rechthoek en $(\Zn{4}, \plusn)$ hebben allebei vier elementen. De symmetriegroep van de driehoek en $S_3$ hebben er allebei zes. Zijn het telkens ``dezelfde'' groepen? Om dat te onderzoeken, kijk je hoe vaak je een element met zichzelf moet samenstellen om bij het neutraal element uit te komen.

\begin{theorie}
In een groep $(G, \ast)$ zijn de \textbf{\textcolor{accent}{machten}} van $a$:

\vspace{1mm}
\qquad $\formule{a^0 = e}, \quad \formule{a^1 = a}, \quad \formule{a^2 = a \ast a}, \quad \formule{a^3 = a \ast a \ast a}, \ \ldots$

\vspace{1mm}
De \textbf{\textcolor{accent}{orde van een element}} $a$ is het kleinste natuurlijk getal $k \geqslant 1$ waarvoor \quad $\formule{a^k = e}$
\end{theorie}

Bij een optelling is $a^k$ gewoon $a$ opgeteld met zichzelf: in $(\Zn{4}, \plusn)$ is $1^4 = 1 \plusn 1 \plusn 1 \plusn 1 = 0$, dus 1 heeft orde 4. Het neutraal element heeft altijd orde 1. In een oneindige groep heeft niet elk element een orde: in $(\Z, +)$ kom je met $1 + 1 + \cdots$ nooit op 0.

\voorbeeld{}{Bepaal de orde van elk element van de symmetriegroep van de gelijkzijdige driehoek.\label{vb:h3-orde}}{%
\uitleg{We stellen elk element met zichzelf samen tot we $e$ krijgen, en lezen dat af in de tabel.}
\begin{stappen}
e: \text{ orde } 1 & Het neutraal element. \\[1ex]
r^2 \neq e, \ r^3 = e: \text{ orde } 3 & Drie keer $120^\circ$ is een volle draai. \\[1ex]
(r^2)^2 = r \neq e, \ (r^2)^3 = e: \text{ orde } 3 & Drie keer $240^\circ$ is twee volle draaien. \\[1ex]
s_1^2 = s_2^2 = s_3^2 = e: \text{ orde } 2 & Twee keer omklappen om dezelfde as.
\end{stappen}}

\subsection{Groepen vergelijken}

Twee groepen hebben \textit{dezelfde structuur} als je de elementen van de ene zo een andere naam kunt geven dat je precies de tabel van de andere krijgt. Wiskundigen noemen zulke groepen \textit{isomorf}. Bij het hernoemen verandert de orde van een element niet. Hebben twee groepen niet evenveel elementen van elke orde, dan hebben ze dus niet dezelfde structuur.

\voorbeeld{}{Hebben de symmetriegroep van de rechthoek en $(\Zn{4}, \plusn)$ dezelfde structuur?\label{vb:h3-structuur}}{%
\uitleg{Beide groepen hebben vier elementen en zijn commutatief. We vergelijken de orde van de elementen.}
\begin{stappen}
\text{rechthoek: } 1, 2, 2, 2 & $e$ heeft orde 1, en $d^2 = h^2 = v^2 = e$. \\[1ex]
\Zn{4}\text{: } 1, 4, 2, 4 & $0$ heeft orde 1, $2 \plusn 2 = 0$, en 1 en 3 hebben orde 4.
\end{stappen}
In $\Zn{4}$ is er een element van orde 4, in de rechthoek niet: de twee groepen hebben niet dezelfde structuur, al hebben ze evenveel elementen.}

\voorbeeld{}{Toon aan dat de symmetriegroep van de gelijkzijdige driehoek dezelfde structuur heeft als $S_3$.\label{vb:h3-driehoeks3}}{%
\uitleg{Nummer de drie plaatsen van de hoeken: 1 boven, 2 linksonder, 3 rechtsonder. Elke symmetrie stuurt elke plaats naar een plaats: dat is een permutatie. Die naam geven we haar.}
\naast{\begin{minipage}{11cm}
\begin{stappenvrij}{6.4cm}{4.2cm}
r = \perm{1 & 2 & 3 \\ 2 & 3 & 1}, \quad r^2 = \perm{1 & 2 & 3 \\ 3 & 1 & 2} & De hoek boven gaat bij $r$ naar linksonder. \\[3ex]
s_1 = \perm{1 & 2 & 3 \\ 1 & 3 & 2}, \quad s_2 = \perm{1 & 2 & 3 \\ 3 & 2 & 1} & Bij $s_1$ blijft plaats 1, bij $s_2$ plaats 2. \\[3ex]
s_3 = \perm{1 & 2 & 3 \\ 2 & 1 & 3}, \quad e = \perm{1 & 2 & 3 \\ 1 & 2 & 3} & Bij $s_3$ blijft plaats 3.
\end{stappenvrij}
\end{minipage}}{\begin{tikzpicture}[scale=.8,font=\small]
  \driehoekassen
  \driehoek{1}{2}{3}
\end{tikzpicture}}
De zes symmetrieën geven zes verschillende permutaties: alle elementen van $S_3$. Een symmetrie na een symmetrie is dezelfde functie als de permutatie na de permutatie. Dus wordt de tabel van de driehoek met de nieuwe namen precies de tabel van $S_3$.}

Bij het vierkant werkt dat niet: het heeft 8 symmetrieën, maar er zijn $4! = 24$ permutaties van de vier hoekpunten. Niet elke permutatie is een symmetrie. Wissel je alleen twee naburige hoekpunten, dan wordt een zijde een diagonaal, en die is langer.

\begin{center}
\renewcommand{\arraystretch}{1.3}
\begin{tabular}{|c|c|c|}
\hline
\rowcolor{accent!14} orde van de groep & groepen uit deze cursus met dezelfde structuur & commutatief \\ \hline
2 & $(\Zn{2}, \plusn)$, \ $\bigl(\{1, -1\}, \cdot\bigr)$ & ja \\ \hline
3 & $(\Zn{3}, \plusn)$, \ de draaiingen van de gelijkzijdige driehoek & ja \\ \hline
4 & $(\Zn{4}, \plusn)$, \ $\bigl(\{1, i, -1, -i\}, \cdot\bigr)$, \ $\bigl(\Zn{5} \setminus \{0\}, \maaln\bigr)$ & ja \\ \hline
4 & de rechthoek, \ $\bigl(\{1, 3, 5, 7\}, \maaln\bigr)$ in $\Zn{8}$ & ja \\ \hline
6 & $(\Zn{6}, \plusn)$, \ $\bigl(\Zn{7} \setminus \{0\}, \maaln\bigr)$ & ja \\ \hline
6 & de gelijkzijdige driehoek, \ $S_3$ & nee \\ \hline
\end{tabular}
\end{center}

\begin{werkwijze}
Om te tonen dat twee groepen niet dezelfde structuur hebben, vergelijk je het aantal elementen, of ze commutatief zijn, en hoeveel elementen er van elke orde zijn. Verschilt er één ding, dan is de structuur verschillend. Om te tonen dat ze dezelfde structuur hebben, geef je elk element van de ene groep een naam uit de andere, zodat de tabellen gelijk worden.
\end{werkwijze}

\begin{oefeningen}
\oef{h3-orde}

\oef{h3-structuur}
\end{oefeningen}

%========================================================================================
\opdrachtenkop

\begin{oefeningen}
\oef{h3-o1}

\oef{h3-o5}

\oef{h3-o2}

\oef{h3-o4}

\oef{h3-o3}

\oef{h3-o7}

\oef{h3-o8}

\oef{h3-o6}

\oef{h3-o9}

\oef{h3-o11}

\oef{h3-o12}

\oef{h3-o10}

\oef{h3-o13}

\oef{h3-o14}

\oef{h3-o15}

\oef{h3-o16}
\end{oefeningen}

%----------------------------------------------------------------------------------------
% KOPPEL: verbanden met andere leerstof
%----------------------------------------------------------------------------------------
\needspace{20\baselineskip}
\opdrachtsectie{Koppel}

\lkr{Telproblemen staat in hetzelfde jaar als deze cursus: kijk na of de productregel en de faculteit al gezien zijn. Complexe getallen en Matrices staan in het 5de jaar. De oefening met matrices geeft de draaiings- en spiegelingsmatrix: de leerlingen hoeven ze niet zelf op te stellen.}

\begin{oefeningen}
\oef{h3-koppel-1}

\oef{h3-koppel-2}

\oef{h3-koppel-3}
\end{oefeningen}

%----------------------------------------------------------------------------------------
\opdrachtsectie{Test jezelf}

\begin{oefeningen}
\oef{h3-t1}

\oef{h3-t2}

\oef{h3-t3}

\oef{h3-t4}

\oef{h3-t5}

\oef{h3-t6}
\end{oefeningen}

\herhaalschema{%
  \fv{oef:h3-t1}{vb:h3-samenstellen}{oef:h3-driehoek}%
  \fv{oef:h3-t2}{vb:h3-rechthoek}{oef:h3-figuren}%
  \fv{oef:h3-t3}{vb:h3-perm}{oef:h3-perm}%
  \fv{oef:h3-t4}{vb:h3-perminvers}{oef:h3-perm}%
  \fv{oef:h3-t5}{vb:h3-orde}{oef:h3-orde}%
  \fv{oef:h3-t6}{vb:h3-structuur}{oef:h3-structuur}}

\ifzelfstudie
\section*{Maak je eigen samenvatting}
Sluit dit boek en schrijf in je schrift: wat een symmetrie is, waarom de symmetrieën van een figuur een groep vormen, de zes symmetrieën van de gelijkzijdige driehoek en hoe je ze samenstelt, wat een permutatie is, hoe je permutaties samenstelt en inverteert, hoeveel elementen $S_n$ heeft, wat de orde van een element is, en hoe je nagaat of twee groepen dezelfde structuur hebben. Vergelijk daarna met de kaders.
\fi
